% HPC Metrics Homework: Chinese Translation + Full Solutions
\documentclass[12pt]{article}
\usepackage{amsmath, amssymb, geometry}
\geometry{a4paper, margin=1in}
\usepackage{xeCJK}
\setCJKmainfont{SimSun}

\title{高性能计算性能度量题目（中文题目 + 中文解答）}
\author{}
\date{}

\begin{document}
\maketitle

\section{Amdahl 定律}

以下问题互相独立：

\subsection{题 1：顺序部分为 4\% 时 2$^n$ 处理器的加速比与效率}

已知顺序比例：

\[
seq = 0.04
\]

Amdahl 加速比：
\[
S(p) = \frac{1}{seq + \frac{1-seq}{p}}
\]

计算 $p = 2^n, n=2\dots 7$。

结果：

\begin{center}
\begin{tabular}{c|c|c|c}
$n$ & $p$ & 加速比 $S(p)$ & 效率 $E=S(p)/p$ \\ \hline
2 & 4 & 3.57 & 0.89 \\
3 & 8 & 6.41 & 0.80 \\
4 & 16 & 10.03 & 0.63 \\
5 & 32 & 14.08 & 0.44 \\
6 & 64 & 17.27 & 0.27 \\
7 & 128 & 19.42 & 0.15 \\
\end{tabular}
\end{center}

最大加速比：
\[
S_{\max} = \frac{1}{0.04} = 25
\]

\subsection{题 2：由加速比反推顺序比例}

Karp--Flatt 因子：

\[
e = \frac{1/Acc - 1/p}{1 - 1/p}
\]

此 $e$ 表示有效顺序部分（包含通信、同步等开销）。

程序 A 和 B 的结果如下：

\paragraph{程序 A：}
\begin{center}
\begin{tabular}{c|c|c}
p & 加速比 & 顺序比例 $e$ \\ \hline
2 & 1.9 & 0.0526 \\
3 & 2.73 & 0.0496 \\
4 & 3.47 & 0.0494 \\
5 & 4.16 & 0.0485 \\
6 & 4.8 & 0.0479 \\
7 & 5.38 & 0.0477 \\
8 & 5.93 & 0.0478 \\
\end{tabular}
\end{center}

程序 A 的顺序部分几乎恒定，扩展性好。

\paragraph{程序 B：}
\begin{center}
\begin{tabular}{c|c|c}
p & 加速比 & 顺序比例 $e$ \\ \hline
2 & 1.94 & 0.0309 \\
3 & 2.72 & 0.0463 \\
4 & 3.34 & 0.0509 \\
5 & 3.82 & 0.0569 \\
6 & 4.2 & 0.0629 \\
7 & 4.49 & 0.0690 \\
8 & 4.72 & 0.0758 \\
\end{tabular}
\end{center}

程序 B 的 $e$ 随 p 增大，说明并行开销随处理器数量迅速增加，可扩展性差。

\subsection{题 3：一般 Amdahl 模型推导}

\subsubsection*{(a) p 处理器加速比}

\[
S(p) = \frac{1}{seq + \frac{1-seq}{p}}
\]

最大加速比：
\[
S_{\max} = \frac{1}{seq}
\]

\subsubsection*{(b) 处理器数量 $p\to\infty$}

\[
S(\infty)= \frac{1}{seq}
\]

\subsubsection*{(c) 若 seq = k/n}

\[
S(p)=\frac{1}{k/n + (1-k/n)/p}
\]

当 $n\to\infty$：

\[
S(p) \to p
\]

表示程序规模越大，可并行比例越高，加速比越接近线性提升。


\section{标量 / 向量模式}

向量化模式速度为 20 倍，假设向量化比例为 $v$，加速比：

\[
S(v)=\frac{1}{(1-v)+v/20}
\]

\subsection{题 (a) 图形略}
使用上述公式即可绘制加速比随 $v$ 的变化曲线。

\subsection{题 (b) 获得加速比 2 所需向量化比例}

\[
2 = \frac{1}{1-v+v/20}
\]
\[
v \approx 52.6\%
\]

\subsection{题 (c) 最大加速比一半的向量化比例}
最大加速比：$20$。半数为 $10$。

\[
10 = \frac{1}{1-v+v/20}
\]

解得：
\[
v = 95\%
\]

\subsection{题 (d)}
向量化比例为 70\%，若无法升级硬件，可通过提升向量化比例获得相同性能，包括：循环展开、SIMD、数据对齐、重写内存访问模式、使用编译器提示等。


\section{MIPS / MFLOPS 性能度量}

\subsection{题 2.1 Whetstone 基准测试}

时钟频率：$16.67\,\text{MHz}$。

\subsubsection*{(a) MIPS}
协处理器：$CPI=10$：
\[
MIPS \approx 1.67
\]

软件模拟：$CPI=6$：
\[
MIPS \approx 2.78
\]

\subsubsection*{(b) 指令总数}
协处理器，$t=1.08$：
\[
N\approx 1.8\times 10^6
\]

软件模拟，$t=13.6$：
\[
N\approx 3.8\times 10^7
\]

\subsubsection*{(c) 每个浮点操作的整数指令数}
\[
\frac{3.8\times10^7}{195578} \approx 190-200
\]

\subsubsection*{(d) MFLOPS}
协处理器：
\[
MFLOPS = \frac{195578}{1.08\times 10^6} \approx 0.18
\]


\subsection{题 2.2 两台系统对比}

\subsubsection*{(a) MIPS 公式}
无协处理器：
\[
MIPS_{no} = \frac{I+YF}{W\times 10^6}
\]

带协处理器：
\[
MIPS_{cop} = \frac{I+F}{B\times 10^6}
\]

\subsubsection*{(b) 求 I}
已知 $F=8\times10^6, Y=50, W=4$：

\[
I=8\times10^7
\]

\subsubsection*{(c) 求 B}
\[
B=1.1\text{秒}
\]

\subsubsection*{(d) MFLOPS}
\[
MFLOPS \approx 7.27
\]

\subsubsection*{(e) 是否应购买协处理器？}
执行时间从 4 秒降至 1.1 秒，MFLOPS 显著提升，即使 MIPS 降低，整体性能明显提升，因此应购买。


\section{Gustafson--Barsis 定律}

\subsection{题 1}
p=8, seq=0.14：
\[
Acc=8 - 0.98 = 7.02
\]

\subsection{题 2}
要求 $Acc(10)\ge 7.02$：
\[
seq \le 33.1\%
\]


\section{Karp--Flatt 指标}

\subsection{Cray Y-MP/8}
表格略，结果：

\begin{center}
\begin{tabular}{c|c|c|c}
p & 加速比 & 效率 & e \\ \hline
2 & 1.95 & 0.98 & 0.023 \\
3 & 2.88 & 0.96 & 0.021 \\
4 & 3.76 & 0.94 & 0.021 \\
8 & 6.96 & 0.87 & 0.021 \\
\end{tabular}
\end{center}

扩展性极佳。

\subsection{Aliant FX/40}

\begin{center}
\begin{tabular}{c|c|c|c}
p & 加速比 & 效率 & e \\ \hline
2 & 1.90 & 0.95 & 0.053 \\
3 & 2.65 & 0.88 & 0.065 \\
4 & 3.22 & 0.81 & 0.080 \\
\end{tabular}
\end{center}

扩展性较差，开销随 p 增大明显。

\end{document}


